\textbf{Phần 1. Trắc nghiệm lựa chọn}
\setcounterpageref{ex}{0}

\begin{ex} %[3D2V15]
Tính giá trị của biểu thức: $M = 7 \times 9 - 48 : 8 + 14$.
\choice
{$68$}
{$73$}
{\True $71$}
{$69$}
\loigiai{
Thực hiện nhân chia trước: $M = 63 - 6 + 14 = 57 + 14 = 71$.
\textbf{Đáp án của câu hỏi là C}
}
\end{ex}

\begin{ex} %[3D2V15]
Tìm $x$ biết: $x \times 8 - 25 = 39$.
\choice
{\True $8$}
{$7$}
{$9$}
{$6$}
\loigiai{
Ta có: $x \times 8 = 39 + 25 = 64 \Rightarrow x = 64 : 8 = 8$.
\textbf{Đáp án của câu hỏi là A}
}
\end{ex}

\begin{ex} %[3D2V15]
Một cửa hàng có $72\text{ kg}$ gạo. Ngày thứ nhất bán được $\frac{1}{8}$ số gạo, ngày thứ hai bán được $\frac{1}{9}$ số gạo. Hỏi sau hai ngày cửa hàng còn lại bao nhiêu ki-lô-gam gạo?
\choice
{$50\text{ kg}$}
{$52\text{ kg}$}
{$54\text{ kg}$}
{\True $55\text{ kg}$}
\loigiai{
Ngày 1 bán: $72 : 8 = 9\text{ kg}$.\
Ngày 2 bán: $72 : 9 = 8\text{ kg}$.\
Còn lại: $72 - (9 + 8) = 72 - 17 = 55\text{ kg}$.
\textbf{Đáp án của câu hỏi là D}
}
\end{ex}

\begin{ex} %[3D2V15]
Trong một phép chia cho $7$, thương bằng $8$ và số dư là số dư lớn nhất có thể. Số bị chia là
\choice
{$56$}
{\True $62$}
{$61$}
{$63$}
\loigiai{
Số dư lớn nhất là 6.\
Số bị chia là: $7 \times 8 + 6 = 56 + 6 = 62$.
\textbf{Đáp án của câu hỏi là B}
}
\end{ex}

\begin{ex} %[3D2V15]
Tính nhanh giá trị của biểu thức: $8 \times 4 + 8 \times 5 + 8$.
\choice
{$72$}
{$88$}
{\True $80$}
{$96$}
\loigiai{
Ta viết $8 = 8 \times 1$. Đặt thừa số chung:\
$$8 \times (4 + 5 + 1) = 8 \times 10 = 80$$
\textbf{Đáp án của câu hỏi là C}
}
\end{ex}

\begin{ex} %[3D2V15]
Một bác nông dân nuôi đàn vịt và đàn gà có tất cả $63$ con. Biết số con gà bằng $\frac{1}{7}$ tổng số con. Hỏi bác nông dân nuôi số con vịt nhiều hơn số con gà bao nhiêu con?
\choice
{\True $45$ con}
{$36$ con}
{$48$ con}
{$42$ con}
\loigiai{
Số gà là: $63 : 7 = 9$ (con).\
Số vịt là: $63 - 9 = 54$ (con).\
Vịt hơn gà: $54 - 9 = 45$ (con).
\textbf{Đáp án của câu hỏi là A}
}
\end{ex}

\begin{ex} %[3D2V15]
Tìm $x$ biết: $x : 9 + 34 = 42$.
\choice
{$63$}
{$81$}
{$54$}
{\True $72$}
\loigiai{
Ta có: $x : 9 = 42 - 34 = 8 \Rightarrow x = 8 \times 9 = 72$.
\textbf{Đáp án của câu hỏi là D}
}
\end{ex}

\begin{ex} %[3D2V15]
Điền dấu thích hợp: $6 \times 9 - 18 \dots 7 \times 5 + 2$.
\choice
{$>$}
{\True $<$}
{$=$}
{Không so sánh được}
\loigiai{
Vế trái: $54 - 18 = 36$.\
Vế phải: $35 + 2 = 37$.\
Vì $36 < 37$ nên điền dấu $<$. 
\textbf{Đáp án của câu hỏi là B}
}
\end{ex}

\begin{ex} %[3D2V15]
Một người mua $6$ hộp bút, mỗi hộp có $8$ chiếc bút chì. Người đó đem tặng $\frac{1}{4}$ số bút chì đó cho các bạn học sinh nghèo. Hỏi người đó còn lại bao nhiêu chiếc bút chì?
\choice
{$32$ chiếc bút}
{$40$ chiếc bút}
{\True $36$ chiếc bút}
{$30$ chiếc bút}
\loigiai{
Tổng số bút chì là: $6 \times 8 = 48$ chiếc.\
Số bút đã tặng: $48 : 4 = 12$ chiếc.\
Số bút còn lại: $48 - 12 = 36$ chiếc.
\textbf{Đáp án của câu hỏi là C}
}
\end{ex}

\begin{ex} %[3D2V15]
Có $54$ bạn học sinh xếp đều thành 6 hàng. Nếu muốn xếp số học sinh đó thành các hàng, mỗi hàng có 9 bạn thì xếp được bao nhiêu hàng?
\choice
{\True $6$ hàng}
{$5$ hàng}
{$7$ hàng}
{$8$ hàng}
\loigiai{
Xếp được số hàng là: $54 : 9 = 6$ hàng.
\textbf{Đáp án của câu hỏi là A}
}
\end{ex}

\begin{ex} %[3D2V15]
Cho dãy số: $7; 14; 21; 28; \dots$ Tổng của số hạng thứ 3 và số hạng thứ 7 trong dãy là
\choice
{$63$}
{$77$}
{$68$}
{\True $70$}
\loigiai{
Số thứ 3 là: $7 \times 3 = 21$. Số thứ 7 là: $7 \times 7 = 49$.\
Tổng là: $21 + 49 = 70$.
\textbf{Đáp án của câu hỏi là D}
}
\end{ex}

\begin{ex} %[3D2V15]
Một số tự nhiên chia cho $8$ được thương bằng $6$ và dư $5$. Khi lấy số đó chia cho $9$ thì số dư là
\choice
{$7$}
{\True $8$}
{$6$}
{$5$}
\loigiai{
Số đó là: $8 \times 6 + 5 = 53$.\
Lấy $53 : 9 = 5$ dư 8.
\textbf{Đáp án của câu hỏi là B}
}
\end{ex}

\textbf{Phần 2. Trắc nghiệm đúng sai}
\setcounterpageref{ex}{0}

\begin{ex} %[3D2V15]
Xét tính đúng sai của các nhận định tổng hợp Chủ đề 2:
\choiceTF
{\True Một số chia hết cho cả 2 và 3 thì chia hết cho 6.}
{Một phần tám của 64 lớn hơn một phần bảy của 63.}
{\True Trong phép chia cho 9, số dư lớn nhất bằng 8.}
{\True Tích $7 \times 8$ bằng tích $8 \times 7$.}
\loigiai{
\begin{itemize}
    \item Ý a) Đúng: Dấu hiệu chia hết cơ bản.
    \item Ý b) Sai: $\frac{1}{8}$ của 64 là 8; $\frac{1}{7}$ của 63 là 9. Vì $8 < 9$ nên nhận định Sai.
    \item Ý c) Đúng: Số dư luôn bé hơn số chia.
    \item Ý d) Đúng: Tính chất giao hoán của phép nhân.
\end{itemize}
\begin{itemize}
    \item a) \textbf{Đúng}.
    \item b) \textbf{Sai}.
    \item c) \textbf{Đúng}.
    \item d) \textbf{Đúng}.
\end{itemize}
}
\end{ex}

\begin{ex} %[3D2V15]
Một người thợ làm việc trong 6 ngày, mỗi ngày làm được 8 sản phẩm. Người đó đem bán đi một nửa số sản phẩm:
\choiceTF
{\True Người đó làm được tất cả 48 sản phẩm.}
{\True Người đó đã bán đi 24 sản phẩm.}
{Số sản phẩm còn lại đủ để chia đều vào 5 hộp, mỗi hộp 5 sản phẩm.}
{\True Số sản phẩm còn lại chia vừa đủ vào 4 hộp loại 6 sản phẩm.}
\loigiai{
\begin{itemize}
    \item Ý a) Đúng: $6 \times 8 = 48$ sản phẩm.
    \item Ý b) Đúng: $48 : 2 = 24$ sản phẩm.
    \item Ý c) Sai: Cần $5 \times 5 = 25$ sản phẩm, mà chỉ còn 24 sản phẩm.
    \item Ý d) Đúng: $24 : 4 = 6$ sản phẩm.
\end{itemize}
\begin{itemize}
    \item a) \textbf{Đúng}.
    \item b) \textbf{Đúng}.
    \item c) \textbf{Sai}.
    \item d) \textbf{Đúng}.
\end{itemize}
}
\end{ex}

\textbf{Phần 3. Trả lời ngắn}
\setcounterpageref{ex}{0}

\begin{ex} %[3D2V15]
Tính giá trị của biểu thức: $8 \times 7 - 45 : 9$.
\shortans{51}
\loigiai{
\textbf{Đáp số:} $51$\\
$8 \times 7 - 45 : 9 = 56 - 5 = 51$.
}
\end{ex}

\begin{ex} %[3D2V15]
Tìm một số biết rằng $\frac{1}{7}$ của số đó bằng 8.
\shortans{56}
\loigiai{
\textbf{Đáp số:} $56$\\
$8 \times 7 = 56$.
}
\end{ex}

\begin{ex} %[3D2V15]
Tìm $x$ biết: $x : 6 = 14 - 7$.
\shortans{42}
\loigiai{
\textbf{Đáp số:} $42$\\
$x : 6 = 7 \Rightarrow x = 7 \times 6 = 42$.
}
\end{ex}

\begin{ex} %[3D2V15]
Một lớp có 32 học sinh chia đều thành 4 tổ. Hỏi 3 tổ có bao nhiêu học sinh?
\shortans{24}
\loigiai{
\textbf{Đáp số:} $24$\\
$32 : 4 = 8$; $8 \times 3 = 24$ bạn.
}
\end{ex}

\begin{ex} %[3D2V15]
Tính nhanh: $9 \times 2 + 9 \times 3 + 9 \times 5$.
\shortans{90}
\loigiai{
\textbf{Đáp số:} $90$\\
$9 \times (2 + 3 + 5) = 9 \times 10 = 90$.
}
\end{ex}

\begin{ex} %[3D2V15]
Một hình tam giác có độ dài ba cạnh đều bằng 8 cm. Chu vi hình tam giác đó là bao nhiêu xăng-ti-mét?
\shortans{24}
\loigiai{
\textbf{Đáp số:} $24$\\
Chu vi là: $8 \times 3 = 24\text{ cm}$.
}
\end{ex}
